% unidades/07-pasiva-perjuicio.tex
\chapter{🌧 \ruby{迷惑}{めいわく}の\ruby{受け身}{うけみ} — Pasiva de Perjuicio}
\unidadheader{07}{迷惑}{迷惑の受け身}{Expresar que uno fue negativamente afectado por la acción de otro}

\begin{tcolorbox}[vocabbox, title={\emoji{pushpin} Vocabulario}]
\begin{tabularx}{\linewidth}{llX}
\toprule
\textbf{Japonés} & \textbf{Español} & \textbf{Tipo} \\ \midrule
\vocabitem{降る}{ふる}{caer (lluvia/nieve)}{v. G1}
\vocabitem{泣く}{なく}{llorar}{v. G1}
\vocabitem{騒ぐ}{さわぐ}{hacer ruido / alborotar}{v. G1}
\vocabitem{邪魔する}{じゃまする}{interrumpir / molestar}{v. suru}
\vocabitem{死ぬ}{しぬ}{morir}{v. G1}
\vocabitem{逃げる}{にげる}{escapar / huir}{v. G2}
\vocabitem{泥棒}{どろぼう}{ladrón}{sustantivo}
\bottomrule
\end{tabularx}
\end{tcolorbox}

\subsection*{\emoji{speech-balloon} Frases Clave}
\frase{\ruby{雨}{あめ}に\ruby{降}{ふ}られた。}{Me cayó la lluvia encima. (fui perjudicado por la lluvia)}
\frase{\ruby{隣}{となり}の\ruby{人}{ひと}に\ruby{騒}{さわ}がれて\ruby{眠}{ねむ}れなかった。}{No pude dormir porque el vecino hizo ruido.}
\frase{\ruby{泥棒}{どろぼう}に\ruby{自転車}{じてんしゃ}を\ruby{盗}{ぬす}まれました。}{Me robaron la bicicleta.}
\frase{\ruby{子供}{こども}に\ruby{泣}{な}かれて\ruby{困}{こま}った。}{Me causó problema que el niño llorara.}

\subsection*{\emoji{gear} Gramática}
\patron{Pasiva de perjuicio con V intransitivo}
  {[主語]は [原因]に [V-intrans.-pasiva]}
  {\ej{\ruby{雨}{あめ}に\ruby{降}{ふ}られた。}{La lluvia me cayó encima — fui perjudicado por la lluvia.}
   \ej{\ruby{友達}{ともだち}に\ruby{来}{こ}られて、\ruby{勉強}{べんきょう}できなかった。}{Mi amigo vino (sin avisar) y no pude estudiar.}
   \ej{\ruby{赤}{あか}ちゃんに\ruby{泣}{な}かれて、\ruby{困}{こま}った。}{Me causó problemas que el bebé llorara.}}

\patron{Pasiva de perjuicio con pérdida de posesión}
  {[主語]は [人]に [もの]を [V-pasiva]}
  {\ej{\ruby{泥棒}{どろぼう}に\ruby{財布}{さいふ}を\ruby{盗}{ぬす}まれました。}{El ladrón me robó la billetera.}
   \ej{\ruby{弟}{おとうと}に\ruby{ゲーム}{ゲーム}を\ruby{壊}{こわ}された。}{Mi hermano menor me rompió el videojuego.}
   \ej{\ruby{猫}{ねこ}にさかなを\ruby{食}{た}べられてしまった。}{El gato se comió mi pescado.}}

\comparacion{Pasiva directa vs Pasiva de perjuicio}
  {Pasiva directa\\\small evento neutral o positivo\\\small \ruby{先生}{せんせい}に\ruby{褒}{ほ}められた\\\textit{Fui elogiado (positivo)}}
  {Pasiva de perjuicio\\\small evento negativo para el sujeto\\\small \ruby{雨}{あめ}に\ruby{降}{ふ}られた\\\textit{Me cayó la lluvia (negativo)}}
  {La diferencia es el tono afectivo. La pasiva de perjuicio expresa que quien habla fue afectado negativamente, incluso con verbos intransitivos.}

\coloquial{〜られて困った}{〜られてまいった / マジ困った}{
  〜られて困った es muy frecuente para expresar incomodidad o queja. En casual: まいった es una exclamación de rendición o frustración.}

\culturabox{\ruby{被害者意識}{ひがいしゃいしき} — La perspectiva de la víctima}{
  El japonés tiene una rica gramática para expresar haber sido afectado negativamente. Esta voz pasiva de perjuicio permite al hablante colocarse como víctima de las circunstancias — incluso de la lluvia. Refleja una sensibilidad cultural hacia el impacto que las acciones ajenas tienen en uno mismo.
}

\lectura{\ruby{電車}{でんしゃ}の\ruby{中}{なか}で\\[0.4em]
  \ruby{昨日}{きのう}の\ruby{朝}{あさ}は\ruby{最悪}{さいあく}だった。まず\ruby{雨}{あめ}に\ruby{降}{ふ}られてぬれてしまった。電車では\ruby{隣}{となり}の\ruby{人}{ひと}に\ruby{足}{あし}を\ruby{踏}{ふ}まれ、うるさい\ruby{人}{ひと}に\ruby{電話}{でんわ}で\ruby{騒}{さわ}がれた。\ruby{帰}{かえ}ってきたら\ruby{猫}{ねこ}に\ruby{夕食}{ゆうしょく}を\ruby{食}{た}べられてしまっていた。}
{En el tren\\[0.4em]
  La mañana de ayer fue terrible. Primero me cayó la lluvia y me mojé. En el tren, me pisaron y un pasajero ruidoso habló por teléfono. Cuando llegué a casa, el gato se había comido mi cena.}

\begin{tcolorbox}[ejerciciobox, title={\emoji{pencil} Ejercicios}]
\begin{enumerate}
  \item Transforma a pasiva de perjuicio: \ruby{雨}{あめ}が\ruby{降}{ふ}った。(yo sufrí)
  \item Construye: \ruby{友達}{ともだち} + \ruby{財布}{さいふ} + \ruby{使}{つか}う (pérdida de posesión)
  \item ¿En qué situaciones usarías la pasiva de perjuicio? Da un ejemplo propio.
\end{enumerate}
\textbf{Respuestas:}
1) \ruby{雨}{あめ}に\ruby{降}{ふ}られた \quad
2) \ruby{友達}{ともだち}に\ruby{財布}{さいふ}を\ruby{使}{つか}われた \quad
3) (libre)
\end{tcolorbox}

\begin{tcolorbox}[kanjibox, title={\emoji{mahjong} Kanjis}]
\begin{tabularx}{\linewidth}{cllXX}
\toprule
\textbf{Kanji} & \textbf{On} & \textbf{Kun} & \textbf{Significado} & \textbf{Vocab} \\ \midrule
\kanjirow{迷}{めい}{\ruby{まよ}{う}}{perderse/dudar}{\ruby{迷惑}{めいわく} / \ruby{迷}{まよ}う}
\kanjirow{惑}{わく}{\ruby{まど}{う}}{confundirse}{\ruby{迷惑}{めいわく} / \ruby{困惑}{こんわく}}
\kanjirow{盗}{とう}{\ruby{ぬす}{む}}{robar}{\ruby{盗}{ぬす}む / \ruby{盗難}{とうなん}}
\kanjirow{騒}{そう}{\ruby{さわ}{ぐ}}{alborotar}{\ruby{騒}{さわ}ぐ / \ruby{騒音}{そうおん}}
\bottomrule
\end{tabularx}
\end{tcolorbox}
